\subsection{When replay supplies a signal that fresh groups lack}
\label{app:theory-gradient-availability}

High accuracy can remove the reward contrasts that fresh-group learning
needs. A stored target can still supply a learning signal in that regime.


The collapse theorem concerns an expected direction. A separate question is
how often a finite group supplies any task signal, and whether a stored mode
must reappear in that group before receiving a nonzero update from the replay objective.

\begin{lemma}[Fresh-group starvation and replay signal]
\label{lem:replay-gradient-availability}
Fix a categorical policy with correctness probability $P$, and draw $G\ge2$
fresh responses independently. For GRPO, Dr.GRPO, and the implemented centered
MaxRL estimator, the probability of a mixed-reward group is
\[
 h_G(P)=1-P^G-(1-P)^G\le G\min\{P,1-P\}.
\]
Every all-equal group has zero fresh task gradient. Thus $h_G$ is the
probability of an opportunity for a nonzero task gradient, rather than a
lower bound on the gradient magnitude.

For a fixed bank with normalized positive weights $w_b$, let
$R_w=-\sum_b w_b\log p_b$ and $\rho_w>0$. Extend $w$ by zero outside the
bank to obtain $\bar w$. The deterministic replay ascent vector is
\[
 v^{\rm replay}=-\rho_w\nabla R_w=\rho_w(\bar w-p).
\]
It vanishes only at $p=\bar w$. In particular, if $p_b\le w_b/2$ for a banked
mode, then $v^{\rm replay}_b\ge\rho_w w_b/2$, even when $b$ is absent from the
fresh group. For a uniform bank of $k\ge2$ modes, as $p\to\mathbf e_b$ at
one banked correct mode, $h_G(P)\to0$ whereas
$\|v^{\rm replay}\|_2\to\rho_w\sqrt{1-1/k}>0$.
The replay norm therefore stays positive as fresh task-signal opportunities vanish.
\end{lemma}

\begin{proof}
The all-correct and all-incorrect events are disjoint and have probabilities
$P^G$ and $(1-P)^G$. A mixed group requires both a correct and an incorrect
sample; the union bound gives the two bounds $GP$ and $G(1-P)$.
For $R_w=-\sum_b w_b\log p_b$, the softmax score identity gives
$\partial R_w/\partial z_a=p_a-\bar w_a$, proving the replay formula and
coordinate bound. Finally,
\[
 \|\bar w-\mathbf e_b\|_2^2=(1-1/k)^2+(k-1)/k^2=1-1/k,
\]
which gives the strictly positive limiting norm for $k\ge2$.
\end{proof}

\paragraph{Worked example: a rare key need not reappear to receive replay.}
Consider the hypothetical policy $p=(0.600,0.389,0.001,0.010)$, with the
first three categories correct, and take $G=16$. Its $P=0.99$ gives
\[
 h_{16}(0.99)=1-0.99^{16}-0.01^{16}\approx0.14854.
\]
Only about $14.85\%$ of groups offer a fresh task-gradient opportunity.
The third key is absent with probability $0.999^{16}\approx0.98412$.
Nevertheless, uniform replay over the three correct keys with $\rho_w=0.1$
has third logit coordinate
$v^{\rm replay}_3=0.1(1/3-0.001)\approx0.03323$, including on groups that
miss that key. The two sampling probabilities describe different events;
the replay value is the deterministic full-bank vector from the lemma.

Here the replay vector averages the fixed bank; a stochastic draw of one
exemplar has this vector only in expectation. Its nonzero logit coordinate
does not guarantee a nonzero neural-parameter gradient after multiplication
by a possibly degenerate Jacobian. Nor does replay necessarily increase the
next mixed-group probability: $h_G$ is largest at $P=1/2$, so increasing an
already high $P$ decreases it. The precise benefit is a separate verified
learning signal when the fresh task gradient is unavailable.

